% ch11_d §11.10–11.11 (书页 405–414 / p419–p428)
%--C-TAIL--
%JOIN%路径必须为零。序参量 $\Psi(\mathbf{r})$ 本身必须是超导体内位置的单值函数，因此当我们绕回到出发点时，其相位角必须回到同一数值，再加上 $2\pi$ 的某个整数倍。由初等矢量演算可得
\begin{equation*}
\Phi = \oint \mathbf{A}\cdot\mathbf{dl} = (\hbar c/e^{*}) \oint \nabla\chi(\mathbf{r})\cdot\mathbf{dl} = n\,2\pi\hbar c/e^{*},
\tag{11.83}
\end{equation*}
其中 $n$ 为整数。\emph{全磁通}（fluxoid）$\Phi$——即穿过回路的总磁通——以 $2\pi\hbar c/e^{*}$ 为单位而量子化。

\emph{磁通量子化}（flux quantization）与液 $^{4}$He 中类似的\emph{环流量子化}（quantization of circulation）一道，是量子力学最引人注目的宏观表现之一。实验发现，冻结在超导系统中的磁通以下列量子为单位跳变
\begin{equation*}
\Phi_0 = 2\pi\hbar c/2e = 2.07\times 10^{-7}\,\mathrm{G\,cm^{2}}.
\tag{11.84}
\end{equation*}
超流粒子携带的电荷恰好是一个电子对的电荷，即
\begin{equation*}
e^{*} = 2e.
\tag{11.85}
\end{equation*}
除了这一重整化之外，基本的 London 假设由此得到完全的实验证实。

对电荷 (11.85) 的确正是我们根据 BCS 理论所应预期的结果。要从第一性原理推导诸如 (11.80) 那样的方程（同时顾及“局域性”近似），需要依靠 \emph{Gor'kov 方程}，它以 Green 函数的语言写成。在这一形式体系中，序参量 $\Psi(\mathbf{r})$ 联系着一个“反常传播子”（anomalous propagator）
\begin{equation*}
F_{\uparrow\downarrow}(\mathbf{r}, t; \mathbf{r}', t') = \langle \psi_\uparrow(\mathbf{r}, t)\,\psi_\downarrow(\mathbf{r}', t') \rangle
\tag{11.86}
\end{equation*}
它度量的是自旋相反的一对电子的波函数在 $\mathbf{r}$ 与 $\mathbf{r}'$ 两点、$t$ 与 $t'$ 两时刻之间的平均相位相干性。正如我们在 (11.76) 中看到的，这个量不消失正是超流态的特征；而在金属的“正常”态中，它本来会自动消失。

\section{超导结}

超导体中的能隙可以通过 Giaever 隧穿（Giaever tunnelling）直接观测。如 §6.8 中那样，我们研究夹在譬如正常金属与超导体之间的一层薄绝缘膜的电流/电压特性。正常电子只有当有空着的准粒子态来接纳它们时才能进入超导体。在远低于 $T_c$ 的温度下，电流必须为零，直到外加偏压 $\phi$ 达到能隙 $\Delta$（§11.4），然后急剧跳升（图 210）。这种结的透射系数 $\mathcal{T}$ 难以计算，但电流随电压的变化使我们得以细致检验能隙区中准粒子态的态密度。

%FIG{210}: {准粒子从正常金属隧穿进入超导体：(a) 零偏压时；(b) $\phi > \Delta$ 时；(c) 电流/电压特性，显示温度的影响。}

实际上，由于 §6.8 中已经讨论过的原因，这一转变并不完全陡峭。对于强耦合超导体，即 (11.22) 中参量 $\mathcal{N}(\mathcal{E}_F)V$ 很大的情形，电子–声子相互作用 (11.5) 足以在准粒子中产生相当比例的\emph{声子协助隧穿}（phonon-assisted tunnelling）（参看图 116）。于是，在隧穿电流中观察到的精细结构可以给出关于超导体声子谱的详细信息。

更引人注目的现象是\emph{相干隧穿}（coherent tunnelling），此时有超流电流流过结。它不仅可以通过两个超导体之间的薄绝缘膜来观测，还可以透过正常金属中相当厚的一层（5000 Å）来观测——借助邻近效应（proximity effect），这层正常金属可以暂时变成超导的。换句话说，真正超导体中配对相互作用的强度，足以在电子穿过势垒时保持其间的相位关联，并在另一侧重新建立起相干性。

对于所有这类现象，§11.9 的宏观序参量形式体系显然是合用的。这个函数在这样的边界上如何行为？结左方的 $\Psi_L$ 与右方的 $\Psi_R$ 之间有什么关系？

既然穿过势垒的渗漏是一个随时间变化的过程，我们就需要考虑 $\Psi$ 的运动方程。按正则量子力学的说法，这涉及一个超流粒子的能量，我们把它取为某个值 $\mu$。于是我们为序参量写出一个含时 Schrödinger 方程，即
\begin{equation*}
i\hbar\frac{\partial\Psi}{\partial t} = \mu\Psi.
\tag{11.87}
\end{equation*}
当然，在普通的超流电流中，$\mu$ 应当是超流配对的化学势——即在 (11.27) 中用作能量原点的量 $\zeta$ 的两倍——因而在整个回路上严格处处相同。这样的系统因此在时间上表现为定态，正如 (11.80) 诸式所隐含的那样。

然而，对于通过两个相同超导体之间的结的透射，我们必须考虑耦合方程
\begin{equation*}
\left.\begin{array}{l}
i\hbar\dfrac{\partial\Psi_L}{\partial t} = \mu_L\Psi_L + \mathcal{T}\Psi_R,\\[8pt]
i\hbar\dfrac{\partial\Psi_R}{\partial t} = \mu_R\Psi_R + \mathcal{T}\Psi_L,
\end{array}\right\}
\tag{11.88}
\end{equation*}
其中我们容许超流序以速率 $\mathcal{T}$ 双向转移。注意，这个系数在这里以\emph{矩阵元}的身份出现，虽然它起初定义为准粒子隧穿的\emph{跃迁概率}。这是因为超流由\emph{电子对}组成，电子对的隧穿速率预期正比于单电子速率的\emph{平方}。

方程 (11.88) 具有一般的含时解，其中超流密度 $|\Psi_L|^{2}$ 与 $|\Psi_R|^{2}$ 以相等而相反的速率变化，仿佛超流粒子正从结的一侧被输运到另一侧。我们应当把这个解解释为一个（单位面积的）电流密度
\begin{equation*}
J_s = e^{*}\frac{\partial}{\partial t}|\Psi_R|^{2} = \frac{2}{\hbar}\,\mathcal{T}\,n_s e^{*}\sin(\chi_R - \chi_L),
\tag{11.89}
\end{equation*}
其中 $\Psi_R$ 与 $\Psi_L$ 之间的相位差（参看 (11.81)）随时间按下列方式变化
\begin{equation*}
\frac{\partial}{\partial t}(\chi_R - \chi_L) = -\frac{1}{\hbar}(\mu_R - \mu_L).
\tag{11.90}
\end{equation*}
当然，这些方程可以在 Gor'kov 形式体系中从 BCS 理论严格推导出来，只需用一个隧穿哈密顿量来表示电子从一个区域向另一个区域的转移。这里注意，超流电流 (11.89) 与准粒子隧穿系数 $\mathcal{T}$ 成正比，这表明相位相干性再一次发挥了作用。

现在假设我们在结两端加上一个小的恒定电压 $\phi$（小于能隙宽度，以避免准粒子激发）。这相当于超流粒子的化学势有了 $2e\phi$ 的差。方程 (11.89) 与 (11.90) 描写一个随时间振荡的超流电流，其频率为 $2e\phi/\hbar$。这就是交流 Josephson 效应（A.C. Josephson effect），它再次为量子现象中的相位相干性提供了直接的宏观证据，并且已被证明是测量重要原子常数之比 $e/\hbar$ 的极为精确的方法。

在磁场存在时，这些公式需要按式 (10.3) 作通常的修改，以保持矢量势定义中的规范不变性。一个典型的情形如下。假设两个超导体之间有一个面积较大的结（图 211）。由于绝缘势垒本身不是超导的，磁通可以穿过结区，方向平行于分隔面。今考虑成对的点，如 $L_1,R_1$；$L_2,R_2$，位于结的两侧，深入到各自内部足以避开穿透区，但在平行于表面的方向上相距颇远。假设已知相位差恰为
\begin{equation*}
\Delta\chi_1 = \chi(R_1) - \chi(L_1)
\tag{11.91}
\end{equation*}
——它跨接 $L_1R_1$；那么跨接 $L_2R_2$ 的相应相位差 $\Delta\chi_2$ 是多少？

考虑一条从 $L_2$ 到 $L_1$ 的路径，沿这条路径我们对式 (11.82) 的两边积分。正如推导 (11.83) 时那样，$\mathbf{v}_s$ 没有贡献——除穿透区之外，它不可能有平行于结面的分量：于是得到
\begin{equation*}
\chi(L_1) - \chi(L_2) = \frac{e^{*}}{\hbar c}\int_{L_2}^{L_1}\mathbf{A}(\mathbf{r})\cdot\mathbf{dl}.
\tag{11.92}
\end{equation*}

%FIG{211}: {Josephson 结。}

类似地，在右边的材料中沿返回的路径，
\begin{equation*}
\chi(R_2) - \chi(R_1) = \frac{e^{*}}{\hbar c}\int_{R_1}^{R_2}\mathbf{A}(\mathbf{r})\cdot\mathbf{dl}.
\tag{11.93}
\end{equation*}
把这两个方程相加，并略去穿过结时右方路径上的平庸贡献，我们得到
\begin{align*}
\Delta\chi_2 &= \Delta\chi_1 + \frac{e^{*}}{\hbar c}\oint \mathbf{A}\cdot\mathbf{dl}\\
&= \Delta\chi_1 + 2\pi\Phi(1,2)/\Phi_0,
\tag{11.94}
\end{align*}
其中 $\Phi_0$ 是量子化磁通的单位，即 (11.84)，而 $\Phi(1,2)$ 是穿过势垒、介于直线 $L_1R_1$ 与直线 $L_2R_2$ 之间的磁通。换言之，当我们在结区内四处移动时，序参量跨边界的相对相位差是变化的。

于是，宽结上的总电流是外加磁场的复杂函数。我们必须对每一小块面积计算 $(\chi_R - \chi_L)$，由 (11.89) 算出对电流的贡献，然后再对整个结面求积分。例如，对于宽度为 $W$、被总磁通 $\Phi$ 穿过的结，可以导出最大 Josephson 电流的公式，形式为
\begin{equation*}
J \propto \left|\sin(\pi\Phi/\Phi_0)\big/(\pi\Phi/\Phi_0)\right|,
\tag{11.95}
\end{equation*}
这一电流可以在零外加电压下流动，但当改变磁场时它迅速振荡。这种直流 Josephson 效应（D.C. Josephson effect）在双结并联这一更复杂的构型中，提供了测量磁场的极其灵敏的器件。

\section{第二类材料}

对 Meissner 效应理论（§11.7）作细致的分析会揭示一个佯谬。我们曾假设，只要磁场能量一超过超导态与正常态的自由能之差，整块样品就会停止超导——亦即在式 (11.63) 给出的临界场 $H_c$ 处。但这个假设忽略了构造一种非均匀相的可能性：磁通由弥散在超导材料中的正常金属细丝层来承载。这样的\emph{中间态}（intermediate state）显然直到高得多的磁场仍可保持热力学稳定，并且在宏观尺度上仍然是电学超导的。

实验上，在\emph{第二类材料}（type II materials）中观察到的正是与此相仿的状态：Meissner 效应在较低的下临界场 $H_{c1}$ 处消失，但材料直到高得多的场 $H_{c2}$ 仍不显示电阻。这类材料用作电磁铁绕组的重要性使人们的注意力集中于这一现象；这一现象在理论上也有重大的意义。

正如铁磁畴理论（§10.3）中那样，我们需要为同一金属中正常区与超导区之间的边界建立一个理论模型。这意味着，(11.77) 或 (11.79) 中的“超流密度”$n_s$ 必须允许随空间变化。我们还需要一些进一步的方程，来确定序参量 $\Psi(\mathbf{r})$ 的振幅作为 $\mathbf{r}$ 的函数的行为。

Ginzburg--Landau 理论的基本假设是：处于磁场 $H$ 中的超导体，其自由能密度（相对于同温度下零场中的正常金属来计量）是如下形式的函数
\begin{equation*}
g_s - g_n = \alpha|\Psi|^{2} + \tfrac{1}{2}\beta|\Psi|^{4} + \frac{1}{8\pi}H^{2} + \frac{1}{2m^{*}}\left|-i\hbar\nabla\Psi - \frac{e^{*}}{c}\mathbf{A}\Psi\right|^{2}.
\tag{11.96}
\end{equation*}
头两项仅仅表达了一个任意的假设：进入超导相所获得的自由能收益是序参量 $\Psi$ 的函数，并且可以按实量 $|\Psi|^{2}$ 的幂展开。任意系数 $\alpha$ 与 $\beta$ 本身依赖于温度；在 $T_c$ 以下、超导相占优的温度范围内，系数 $\alpha$ 必然变为负；于是零磁场中的超流浓度受 $\beta$ 项的限制，取值为
\begin{equation*}
n_s = |\Psi_0|^{2} = -\alpha/\beta.
\tag{11.97}
\end{equation*}
对于遵从 Meissner 效应的第一类材料，我们可以把 $H^{2}$ 项考虑进来，并把大块材料的临界场认定为
\begin{equation*}
H_c^{2}/8\pi = \alpha^{2}/2\beta,
\tag{11.98}
\end{equation*}
亦即 $g_s$ 再也无法降到 $g_n$ 之下的那一点。同样，我们可以把 (11.97) 代入 London 公式 (11.61)，即
\begin{equation*}
\lambda = \left(\frac{m^{*}c^{2}}{4\pi n_s e^{*2}}\right)^{\frac{1}{2}},
\tag{11.99}
\end{equation*}
以便用同一组参量来表达穿透深度，从而可以估计这些参量的温度依赖性。

(11.96) 的最后一项也是唯象的，但它有牢固的正则基础——它正是磁场中粒子动能的典型量子力学表达式。在 London 方程 (11.55) 与 (11.80) 中，这一项代表“超流电流”的动能，它来自 $\Psi$ 相位的空间变化，如 (11.82) 那样。Ginzburg--Landau 理论的精妙之处在于：恰恰是同一项，系数不改，却能同时描述序参量的模发生空间变化所产生的效应——这些效应本应归因于破坏配对相干性所需的能量。

例如，假设我们想从一个超导区——其中 $\Psi$ 由 (11.97) 给出——过渡到一个 $\Psi = 0$ 的正常区，同时使总自由能的代价最小。在没有磁场并且忽略 $\beta$ 项的情况下，序参量的模必须满足偏微分方程
\begin{equation*}
(\hbar^{2}/2m^{*})\nabla^{2}|\Psi| - |\alpha|\Psi = 0,
\tag{11.100}
\end{equation*}
它由自由能积分经变分法得到。这描写的是 $|\Psi|$ 指数式衰减的情形，其特征长度为
\begin{equation*}
\xi = \left\{\hbar^{2}/2m^{*}|\alpha|\right\}^{\frac{1}{2}}.
\tag{11.101}
\end{equation*}
换言之，我们可以把系数 $\alpha$ 与 §11.8 的相干长度联系起来；在典型的“脏”金属中，这个量近似等于电子平均自由程 $\Lambda$。这个描述显然不足以细致地表现 §11.8 中 Pippard 公式所隐含的非局域行为，但在接近 $T_c$ 的温度下，可以用 Gor'kov 方法从第一性原理予以论证。

论证的下一步，是研究超导材料（其中 $|\Psi| = |\Psi_0|$ 且 $H = 0$）与含有临界场 $H_c$ 的正常材料之间的临界边界处自由能最小的位形。$\Psi$ 的一般偏微分方程要与 $H$ 的 Maxwell 方程组相耦合，如 (11.58)–(11.60) 那样，是相当繁复的，但在两种极端情形下，答案的形式不难求出。

首先假设相干长度 $\xi$ 远大于穿透深度 $\lambda$。由图 212($a$) 可以清楚看到，磁场被从厚度为 $\xi$ 的一层中排除出去，这需要能量
\begin{equation*}
\sigma_{\mathrm{I}} \sim \xi(H_c^{2}/8\pi)
\tag{11.102}
\end{equation*}
（每单位边界面积），而在这一层的大部分区域内，并没有“有序化”自由能的收益作为补偿。表面能近似由 (11.102) 给出，是正的。这是第一类超导电性的典型情形：中间态因需要在正常区与超导区之间制造边界而付出能量，故受到抑制。

另一方面，如果相干长度 $\xi$ 远小于穿透深度 $\lambda$，那么在磁场被排除的区域的大部分范围内，都可以获得有序相的负能量（图 212($b$)）。因此粗略地说，表面能的形式为
\begin{equation*}
\sigma_{\mathrm{II}} \sim -\lambda(H_c^{2}/8\pi).
\tag{11.103}
\end{equation*}
更细致的分析表明，表面能在下述情况下由正变负：
\begin{equation*}
\lambda > \xi/\sqrt{2}.
\tag{11.104}
\end{equation*}
这就是第二类超导电性的基本条件；化学杂质与晶体缺陷显然有利于这一条件，因为它们使相干长度 $\xi$ 缩短，而对 $\lambda$ 等其他宏观超导参量没有明显的影响。

%FIG{212}: {正常区与超导区之间的边界。(a) 第一类，$\xi \gg \lambda$；(b) 第二类，$\xi \ll \lambda$。}

正常相与超导相之间的负表面能显然有利于第二类材料在 $H_{c1}$ 与 $H_{c2}$ 之间的磁场范围内所采取的\emph{混合态}（mixed state）。样品被磁通细丝所穿透，每一根细丝都具有单位全磁通强度，即式 (11.84)。细丝内有一个半径为 $\xi$ 的正常金属圆柱形芯区，外围是半径为 $\lambda$ 的圆柱区域，超流电流必须围绕它流动，以限制磁场的穿透。因此，每根细丝的行为都像超流体中的一根涡线（vortex line），并且倾向于排斥其近邻。当我们增大磁场时，磁通线的密度增大，细丝“结晶”成二维点阵。但间隙材料在宏观尺度上仍保持超导，直到上临界场 $H_{c2}$——此时正常金属的芯区被迫相互接触，整块样品回归正常态。

%FIG{213}: {第二类超导体中磁通细丝的阵列，显示磁力线集中于圆柱形正常材料芯的内部及其周围。超流涡旋在穿透区内绕细丝流动。}
